% Preserved full version of former Section 5.4. Input after \appendix.
% Requires amsmath, booktabs, graphicx, and the generated figure directory.
\providecommand{\AFCaseStudyFigures}{output/pdf/assisted-r9-case-study}

\section{Interventional discrimination: detailed case study}
\label{app:interventional_discrimination}

Similar trajectory errors can conceal different responses to a changed
physical condition. We illustrate this on canonical, obfuscated T1-easy using
three historical models: an assisted Autoformalism Brief-only endpoint
(\emph{R9}), a no-specification endpoint, and a GPT-5.6 Sol model. R9 underwent
assisted sign repair, training-only refitting, and validation-selected pruning;
it is not an unassisted Full run. We compare fixed endpoints from different
seeds and search histories, separately from the matched ablation study above.
Their pooled training NMSEs are similar: 0.0416, 0.0436, and 0.0368,
respectively. Figure~\ref{fig:exp_interventional_case} shows single- and
double-meal training examples; all three models have
NMSE below 0.044 on both. The Sol endpoint is repetition 0, which has the
lowest original validation NMSE among its three repetitions
(0.0142 versus 0.0453 and 0.0310), rather than the best result on the new probes.

\paragraph{An omitted dependency makes a response impossible.}
We change initial tissue glucose $G_t(0)$ by $\pm20\%$, keeping initial plasma
glucose $G_p(0)$, the meal schedule, and all other physical initial states fixed.
This changes initial total glucose mass. The physical simulator regenerates
the target and every permitted auxiliary trajectory together; learned models
retain their fitted parameters and causal initializers. As required by T1-easy,
prediction is conditional on the supplied auxiliary trajectories, not an
autonomous simulation of all physiology. The no-specification model has three
latent states but ignores all physiological auxiliary channels, so
\begin{equation}
  \widehat G_p(t)=F_\theta\bigl(t,u_{[0,t]},G_p(0)\bigr).
  \label{eq:exp_missing_dependency}
\end{equation}
Its fitted latent initial values are constants. Consequently, changing $G_t(0)$
while preserving the arguments of $F_\theta$ leaves its prediction unchanged.
R9 instead retains an active tissue-return term,
\begin{equation}
  \dot G_p = h_\theta(B,H,U_{\rm ii},{\rm EGP},u)
       +0.1377G_t-0.09523G_p,
  \label{eq:exp_active_dependency}
\end{equation}
where $B,H$ are its evolving meal-related contributions. This has the positive
direction of the canonical tissue-return flux $+0.079G_t$, although R9's
coefficients and other fluxes remain approximate. Sol 0 has no additional
latent state and compensating fitted auxiliary gains, including
$+0.9735G_t-18.9703\,\mathrm{EGP}-0.5936G_p$ in its plasma equation.
Its excessive response here follows from these fitted gains; absence of extra
latent states alone does not establish the failure.

\begin{figure*}[t]
  \centering
  \includegraphics[width=\textwidth]{\AFCaseStudyFigures/r9_case_study.pdf}
  \caption{\textbf{Comparable training fit, different intervention responses.}
  Canonical obfuscated T1-easy; symbols are decoded for interpretation.
  (a,d) Original training trajectories: a 60 g meal at minute 0 and meals of
  30/60 g at minutes 0/120. (b,c) Absolute fasting trajectories after changing
  initial tissue glucose by $-20\%$ or $+20\%$, with the same initial plasma
  glucose. (e,f) Each trajectory minus that model's own unperturbed fasting
  control. No-specification cannot respond to the omitted tissue channel;
  validation-selected Sol 0 overreacts; assisted R9 captures more of the
  response but retains baseline error. Parameters and initializers are frozen.
  Both intervention directions are shown; R9 does not uniformly win absolute
  NMSE. Complete trajectory grids and all Sol repetitions accompany the figure.}
  \label{fig:exp_interventional_case}
\end{figure*}

\paragraph{Response fidelity and absolute accuracy answer different questions.}
For intervention $c$ and its matched control $c_0$, define
$\Delta_c G_p=G_p^{(c)}-G_p^{(c_0)}$ and subtract each model's own control to
obtain $\Delta_c\widehat G_p$. The response error is
\begin{equation}
  R_c=\frac{\sum_i[\Delta_c\widehat G_p(t_i)-\Delta_cG_p(t_i)]^2}
                 {\sum_i[\Delta_cG_p(t_i)]^2}.
  \label{eq:exp_response_error}
\end{equation}
All reported contrasts have nonzero denominators; predicting no effect gives
$R_c=1$. For fasting $G_t(0)-20\%$, R9 and Sol 0 have absolute NMSE
0.00932 and 0.01985, and response errors 0.126 and 0.901. No-specification's
response error is exactly 1. The positive perturbation illustrates why both
metrics matter: R9 has a smaller response error than Sol 0 (0.177 versus
1.543), but a larger absolute NMSE (0.01944 versus 0.01178).

We retain the entire five-contrast suite: both tissue perturbations under
fasting and under a 60 g meal at minute 60, plus splitting that meal into
30 g at minutes 60 and 90 versus the unsplit control. Table~\ref{tab:exp_response_robustness}
reports the maximum error within this finite suite for every Sol repetition.
R9's worst observed response error is lower than each Sol repetition's, while
each Sol repetition has lower worst absolute trajectory NMSE. Thus the result
concerns intervention-response fidelity, not uniform predictive superiority or
a worst-case guarantee beyond these probes. On the meal/$G_t(0)-20\%$ probe,
R9's absolute NMSE is 0.0566, versus 0.397--0.605 for D3,
0.728--1.190 for PySR, and 1.190 for SINDy across their three repetitions;
those baselines already have substantially worse training fits.

\begin{table}[t]
  \centering
  \small
  \setlength{\tabcolsep}{5pt}
  \caption{\textbf{All frozen Sol repetitions, without selecting on probe outcomes.}
  Training NMSE uses all 16 trajectories. Maxima use the same five exploratory
  interventions; response errors compare against the relevant matched control.
  Absolute NMSE uses the original training variance. R9 is assisted and these
  historical endpoints do not have matched discovery budgets.}
  \label{tab:exp_response_robustness}
  \begin{tabular}{lccc}
    \toprule
    Model & Train NMSE & Max. NMSE & Max. $R_c$ \\
    \midrule
    Assisted R9 & 0.0416 & 0.0693 & \textbf{0.217} \\
    No specification & 0.0436 & 0.1220 & 1.135 \\
    Sol 0 (validation-selected) & 0.0368 & 0.0345 & 1.594 \\
    Sol 1 & 0.0412 & 0.0474 & 0.366 \\
    Sol 2 & 0.1623 & \textbf{0.0186} & 0.536 \\
    \bottomrule
  \end{tabular}
\end{table}

This example motivates specifying relevant dependencies and evaluating their
responses, rather than using trajectory fit alone to infer mechanistic validity.
It does not show that removing the specification caused the omitted pathway:
T1 does not explicitly mandate $G_t$, and the endpoints are not a matched
causal comparison. Nor does it establish complete mechanism recovery by R9.
The models and probes were chosen through exploratory equation inspection;
these results are reported separately from frozen benchmark test estimates.
